\chapter{Modello Dati}

\section{Struttura Gerarchica XML}

I file XML dell'RNA presentano una struttura gerarchica a tre livelli:

\begin{figure}[h]
\centering
\begin{tikzpicture}[
    level 1/.style={sibling distance=5cm, level distance=1.5cm},
    level 2/.style={sibling distance=3cm, level distance=1.5cm},
    every node/.style={rectangle, rounded corners, draw, fill=primaryblue!10, minimum width=2.5cm, minimum height=0.8cm, font=\small}
]
\node {AIUTO}
    child {node {COMPONENTE\_1}
        child {node {STRUM\_1}}
        child {node {STRUM\_2}}
    }
    child {node {COMPONENTE\_2}
        child {node {STRUM\_3}}
    };
\end{tikzpicture}
\caption{Struttura gerarchica dei dati RNA}
\end{figure}

\section{Normalizzazione in Tre Tabelle}

Il parser trasforma la struttura gerarchica in tre tabelle relazionali normalizzate:

\subsection{Tabella AIUTI}

Entità principale contenente le informazioni sul contributo:

\begin{table}[h]
\centering
\small
\begin{tabular}{@{}lll@{}}
\toprule
\textbf{Campo} & \textbf{Tipo} & \textbf{Descrizione} \\
\midrule
CAR & string & Codice Aiuto Regionale (PK logica) \\
COR & string & Codice Operazione Registro \\
TITOLO\_MISURA & string & Denominazione del bando \\
DATA\_CONCESSIONE & string & Data di erogazione \\
DENOMINAZIONE\_BENEFICIARIO & string & Ragione sociale beneficiario \\
CODICE\_FISCALE\_BENEFICIARIO & string & CF/P.IVA \\
REGIONE\_BENEFICIARIO & string & Regione di appartenenza \\
CUP & string & Codice Unico Progetto \\
ANNO & int32 & Partizione temporale \\
FILE\_SOURCE & string & Tracciabilità file sorgente \\
\bottomrule
\end{tabular}
\caption{Schema tabella AIUTI}
\end{table}

\subsection{Tabella COMPONENTI}

Dettaglio delle componenti di aiuto, in relazione 1:N con AIUTI:

\begin{table}[h]
\centering
\small
\begin{tabular}{@{}lll@{}}
\toprule
\textbf{Campo} & \textbf{Tipo} & \textbf{Descrizione} \\
\midrule
ID\_COMPONENTE\_AIUTO & string & Identificativo componente (PK) \\
CAR\_AIUTO & string & FK verso AIUTI.CAR \\
COR\_AIUTO & string & FK verso AIUTI.COR \\
COD\_REGOLAMENTO & string & Regolamento UE applicato \\
SETTORE\_ATTIVITA & string & Codice ATECO \\
ANNO & int32 & Partizione \\
\bottomrule
\end{tabular}
\caption{Schema tabella COMPONENTI}
\end{table}

\subsection{Tabella STRUMENTI}

Strumenti finanziari applicati, in relazione 1:N con COMPONENTI:

\begin{table}[h]
\centering
\small
\begin{tabular}{@{}lll@{}}
\toprule
\textbf{Campo} & \textbf{Tipo} & \textbf{Descrizione} \\
\midrule
ID\_COMPONENTE\_AIUTO & string & FK verso COMPONENTI \\
COD\_STRUMENTO & string & Codice strumento finanziario \\
IMPORTO\_NOMINALE & float64 & Valore nominale contributo \\
ELEMENTO\_DI\_AIUTO & float64 & Equivalente sovvenzione \\
ANNO & int32 & Partizione \\
\bottomrule
\end{tabular}
\caption{Schema tabella STRUMENTI}
\end{table}

\section{Diagramma ER}

\begin{figure}[h]
\centering
\begin{tikzpicture}[
    entity/.style={rectangle, draw, fill=primaryblue!20, minimum width=3cm, minimum height=1.5cm, font=\bfseries},
    relationship/.style={diamond, draw, fill=warnorange!30, minimum width=1.5cm, aspect=2},
    scale=0.9, transform shape
]
    \node[entity] (aiuti) at (0,0) {AIUTI};
    \node[entity] (comp) at (5,0) {COMPONENTI};
    \node[entity] (strum) at (10,0) {STRUMENTI};
    
    \node[relationship] (r1) at (2.5,0) {1:N};
    \node[relationship] (r2) at (7.5,0) {1:N};
    
    \draw[-] (aiuti) -- (r1);
    \draw[-] (r1) -- (comp);
    \draw[-] (comp) -- (r2);
    \draw[-] (r2) -- (strum);
    
    \node[below=0.3cm of aiuti, font=\small\ttfamily] {CAR, COR};
    \node[below=0.3cm of comp, font=\small\ttfamily] {ID\_COMPONENTE};
    \node[below=0.3cm of strum, font=\small\ttfamily] {IMPORTO, ELEMENTO};
\end{tikzpicture}
\caption{Diagramma ER delle tabelle normalizzate}
\end{figure}

\section{Partizionamento Hive-Style}

L'output è organizzato con \textbf{Hive-style partitioning} sulla colonna \texttt{ANNO}:

\begin{lstlisting}[language=bash, caption=Struttura directory output]
public/parquet/
    aiuti/
        ANNO=2022/
            part-0.parquet
            part-1.parquet
        ANNO=2023/
            part-0.parquet
    componenti/
        ANNO=2022/
        ANNO=2023/
    strumenti/
        ANNO=2022/
        ANNO=2023/
\end{lstlisting}

Questo layout consente:
\begin{itemize}
    \item \textbf{Partition pruning}: query filtrate per anno leggono solo le partizioni rilevanti
    \item \textbf{Parallelismo}: file distinti possono essere letti concorrentemente
    \item \textbf{Gestibilità}: facile eliminazione o aggiornamento di singoli anni
\end{itemize}
